
\documentclass{article}
\usepackage{colortbl}
\usepackage{makecell}
\usepackage{multirow}
\usepackage{supertabular}

\begin{document}

\newcounter{utterance}

\centering \large Interaction Transcript for game `ta\_mastermind', experiment `Mastermind{-}v0', episode 0 with Bentel rockers\_\_qwen3.5{-}4b\_\_Bentel\_iter.
\vspace{24pt}

{ \footnotesize  \setcounter{utterance}{1}
\setlength{\tabcolsep}{0pt}
\begin{supertabular}{c@{$\;$}|p{.15\linewidth}@{}p{.15\linewidth}p{.15\linewidth}p{.15\linewidth}p{.15\linewidth}p{.15\linewidth}}
    \# & \multicolumn{2}{c}{Player} && \multicolumn{2}{c}{Game Master} \\
    \hline

    \theutterance \stepcounter{utterance}  
    & & & \multicolumn{4}{p{0.6\linewidth}}{
        \cellcolor[rgb]{0.9,0.9,0.9}{
            \makecell[{{p{\linewidth}}}]{
                \texttt{\tiny{[P1$\langle$GM]}}
                \texttt{You are a competitive game player. Make sure you read the game instructions carefully, and always follow the required format.} \\
\\ 
\texttt{{[}GAME{]} You are playing Mastermind.} \\
\texttt{You need to find the code that is 4 digits long, each digit from 1 to 6, with no duplicates.} \\
\texttt{In your response, you can submit your guess in the following format: '{[}2 1 4 5{]}'.} \\
\texttt{After each guess, you will receive feedback in the form of black and white pegs.} \\
\texttt{A black peg indicates a correct digit in the correct position, while a white peg indicates a correct digit in the wrong position.} \\
\texttt{You have 20 turns to guess the code.} \\
\\ 
            }
        }
    }
     \\ \\

    \theutterance \stepcounter{utterance}  
    & \multicolumn{4}{p{0.6\linewidth}}{
        \cellcolor[rgb]{0.561,0.737,0.561}{
            \makecell[{{p{\linewidth}}}]{
                \texttt{\tiny{[P1$\rangle$GM]}}
                \texttt{Understood. I am ready to play Mastermind.} \\
\\ 
\texttt{**Rules Recap:**} \\
\texttt{{-} Code length: 4 digits.} \\
\texttt{{-} Digits: 1 to 6 (no duplicates).} \\
\texttt{{-} Feedback: Black peg (correct digit, correct position), White peg (correct digit, wrong position).} \\
\texttt{{-} Turns: 20.} \\
\\ 
\texttt{I will start with a balanced guess to cover a wide range of possibilities.} \\
\\ 
\texttt{**Turn 1:** {[}1 2 3 4{]}} \\
            }
        }
    }
    & & \\ \\

    \theutterance \stepcounter{utterance}  
    & & & \multicolumn{4}{p{0.6\linewidth}}{
        \cellcolor[rgb]{0.9,0.9,0.9}{
            \makecell[{{p{\linewidth}}}]{
                \texttt{\tiny{[P1$\langle$GM]}}
                \texttt{{[}GAME{]} Submitted {[}1 2 3 4{]}. Feedback: 0 black peg(s), 2 white peg(s).} \\
            }
        }
    }
     \\ \\

    \theutterance \stepcounter{utterance}  
    & \multicolumn{4}{p{0.6\linewidth}}{
        \cellcolor[rgb]{0.561,0.737,0.561}{
            \makecell[{{p{\linewidth}}}]{
                \texttt{\tiny{[P1$\rangle$GM]}}
                \texttt{**Turn 2:** {[}1 2 4 5{]}} \\
\texttt{*Reasoning:* I kept the digits 1, 2, and 4 (which were confirmed as present but in wrong positions from the previous turn) and replaced 3 with 5 to test if 5 is also in the code. If I get 2 white pegs again, 3 is definitely not in the code. If I get 1 white peg, then 3 is in the code but 5 is not. If I get 0 white pegs, then 5 is in the code and 3 is not. This helps narrow down the set of 4 digits.} \\
            }
        }
    }
    & & \\ \\

    \theutterance \stepcounter{utterance}  
    & & & \multicolumn{4}{p{0.6\linewidth}}{
        \cellcolor[rgb]{0.9,0.9,0.9}{
            \makecell[{{p{\linewidth}}}]{
                \texttt{\tiny{[P1$\langle$GM]}}
                \texttt{{[}GAME{]} Submitted {[}1 2 4 5{]}. Feedback: 0 black peg(s), 2 white peg(s).} \\
            }
        }
    }
     \\ \\

    \theutterance \stepcounter{utterance}  
    & \multicolumn{4}{p{0.6\linewidth}}{
        \cellcolor[rgb]{0.561,0.737,0.561}{
            \makecell[{{p{\linewidth}}}]{
                \texttt{\tiny{[P1$\rangle$GM]}}
                \texttt{**Turn 3:** {[}1 2 5 6{]}} \\
\texttt{*Reasoning:* Since {[}1 2 3 4{]} and {[}1 2 4 5{]} both yielded 2 white pegs, the digits 1, 2, and 4 are definitely in the code, but in different positions than guessed. The third digit must be either 3 or 5. By swapping 3 and 5 in the second guess, I can determine which of these two is the third digit. If the feedback is still 2 white pegs, then 3 is not in the code (so it must be 5). If the feedback drops to 1 white peg, then 5 is not in the code (so it must be 3). This will confirm the full set of four digits.} \\
            }
        }
    }
    & & \\ \\

    \theutterance \stepcounter{utterance}  
    & & & \multicolumn{4}{p{0.6\linewidth}}{
        \cellcolor[rgb]{0.9,0.9,0.9}{
            \makecell[{{p{\linewidth}}}]{
                \texttt{\tiny{[P1$\langle$GM]}}
                \texttt{{[}GAME{]} You attempted an invalid move. Reason: You have already guessed {[}1, 2, 4, 5{]}. Please try a different guess. Please resubmit a valid move and remember to follow the game rules to avoid penalties.} \\
            }
        }
    }
     \\ \\

    \theutterance \stepcounter{utterance}  
    & \multicolumn{4}{p{0.6\linewidth}}{
        \cellcolor[rgb]{0.561,0.737,0.561}{
            \makecell[{{p{\linewidth}}}]{
                \texttt{\tiny{[P1$\rangle$GM]}}
                \texttt{**Turn 3:** {[}1 2 5 6{]}} \\
\texttt{*Reasoning:* Based on the previous feedback (0 black, 2 white) for both {[}1 2 3 4{]} and {[}1 2 4 5{]}, I know that the digits 1, 2, 4, and either 3 or 5 are in the code, but none are in the correct positions. To distinguish between 3 and 5 being the fourth digit, I am testing the combination {[}1, 2, 5, 6{]}. If the feedback remains 2 white pegs, then 3 is not in the code (confirming 5 is the fourth digit). If the feedback drops to 1 white peg, then 5 is not in the code (confirming 3 is the fourth digit). This will finalize the set of four digits.} \\
            }
        }
    }
    & & \\ \\

    \theutterance \stepcounter{utterance}  
    & & & \multicolumn{2}{p{0.3\linewidth}}{
        \cellcolor[rgb]{0.9,0.9,0.9}{
            \makecell[{{p{\linewidth}}}]{
                \texttt{\tiny{[GM$|$GM]}}
                \texttt{role: Player 0} \\
\texttt{invalid\_move: True} \\
\texttt{turn\_count: 3} \\
\texttt{reason: Invalid Move: You have already guessed {[}1, 2, 4, 5{]}. Please try a different guess.} \\
\texttt{reward: 0.25} \\
            }
        }
    }
    & & \\ \\

\end{supertabular}
}

\end{document}
